\section{Limitations}
\label{sec:limitations}

The factorization is exact, but its interpretation in terms of transported defects relies on the first-order reconstruction, whose accuracy is measured in Section~\ref{sec:res_validity}. Beyond that regime, the nonlinear factor $N_T$ absorbs the difference, and the other factors describe the first-order reconstruction rather than the actual error. All quantities are computed in the Euclidean norm of the standardized state. Other metrics (for example, energy norms or task-specific observables) define other singular vectors and other covariance nulls, and could change the efficiencies and the coherence.

The two systems are small enough ($d=20$ and $N=64$) to form every tail propagator explicitly and to compute full singular value decompositions, which costs $\mathcal O(Td^3)$ per window. For high-dimensional surrogates, $S_T$, $L_T$ and $Q_T$ require only Jacobian--vector products, $\sigma_1(\Phi_k)$ can be obtained by power iteration with Jacobian--vector and vector--Jacobian products, and the null models by randomized estimators with a few random vectors or a low-rank $\Sigma$. The full modal accounting of Eq.~\eqref{eq:spectral} would require Krylov approximations of the leading singular subspaces.

The variance shares describe the variability across initial states of a family of surrogates, not a causal decomposition. The surrogates are a residual MLP and a compact one-dimensional FNO trained on single trajectories, and the accuracy of the L96-2000 surrogates is limited by data coverage rather than by optimization (Section~\ref{sec:res_fidelity}). Larger architectures, latent-space models, surrogates trained on many trajectories and surrogates of high-dimensional systems may show different defect structures. The fine-tuning study uses small budgets and a single validation-selected configuration per surrogate family, and its conclusions are restricted to these protocols. In particular, training from scratch with a multistep objective was not tested.

In the fine-tuning study, the 312 evaluation windows are non-overlapping windows of one held-out trajectory and are not necessarily independent, so that the window-bootstrap intervals are conditional on this trajectory and may be too narrow. They do not describe the variability across test trajectories or training protocols.

Finally, the factorization itself is retrospective. The prospective estimate of Section~\ref{sec:res_ensemble} requires several independently trained surrogates, was tested with five members, two systems and moderate dimensions, and does not detect errors shared by all members.
